\section{OpenAI 核心工作：Post-Training RL Infra}

\subsection{初入 OpenAI：大号实验室}

翁家翌是 OpenAI 第 280 号员工（现在已超过 3000 人）。他对 OpenAI 的初印象是"一个大号的实验室"——没有想象中那么有方法论，但有很多 research 直觉很强的人在里面，可以指明方向。

\begin{importantbox}{OpenAI 的"先进生产力"引入}
一个关键转折点是 Barret、Luke 和 Liam 三人从 Google 加入 John Schulman 的 RL team 之后。他们引入了 Google 的先进生产力——一种哲学理念：\textbf{不要想天才的 idea 和天才的算法，先把 infra 打好，让迭代速度从一周三十次提升到一周三百次}。翁家翌展示了一张图：单位时间的迭代次数和成功率成正比——这本质上就是 RL 的做法（trial and error）。
\end{importantbox}

\subsection{ChatGPT 的诞生：意料之外的指数增长}

翁家翌透露了 ChatGPT 发布的内部视角：发布的初衷只是为了收集一些真实用户数据，预期可能有一两万人用，然后五天之后如果没人就关了。结果用户增长曲线是\textbf{指数级}的。ChatGPT 的爆发完全不是计划出来的，而是"一系列半偶然半必然的化学反应"。

\begin{knowledgebox}{ChatGPT 发布的内部预期}
OpenAI 内部对 ChatGPT 的最初预期非常保守：上线后可能有一两万人使用，如果五天后用户量下滑就关掉。没有人预料到会出现指数级增长。这个故事说明，即使是 AI 前沿的顶级团队，也无法准确预测一个产品的市场反响。真正伟大的产品往往不是计划出来的，而是在偶然和必然的交汇中涌现的。
\end{knowledgebox}

\subsection{搭建 Post-Training RL Infrastructure}

翁家翌在 OpenAI 的核心工作是搭建整个 Post-Training 的 RL Infrastructure。这个 infra 处于整个技术栈的最顶端——最面向"客户"（内部 researcher）的一层。

他为自己设定了一个明确的 reward function：\textbf{最大化自己在 OpenAI Blog 上出现名字的次数}。要实现这个目标，做 infra 比做单个 research 更有效，因为 infra 可以 scale up——大家都用你的 infra，每发一个大模型你的名字就得放上去。主持人评价道："你真的很会给自己写 reward 啊。"

翁家翌解释了这个 reward function 背后的逻辑：如果做单个 research 项目，每次只能影响一个发布；但如果做 infra，所有使用这套 infra 训练出来的模型都会在 contributor list 上带你的名字。并且他擅长写 RL Infra，所以这是一个"非常非常适合的机会"。

\begin{importantbox}{Infra 的核心价值：Bug 修复 = 模型质量}
翁家翌提出了一个犀利的观点：\textbf{每家的 infra 都有不同程度的 bug，谁修的 bug 越多，谁的模型训练就越好。}他甚至猜测 LLaMA 打不过 GPT 可能就是因为 LLaMA 的 infra bug 更多。在他看来，AI 前沿竞争的本质不是算法创新，而是\textbf{infra 的正确性和迭代速度}。
\end{importantbox}

\subsection{下一代 Infra：推倒重来}

翁家翌透露 OpenAI 正在\textbf{推倒重来}，重构内部的 RL Infra。原因是之前那代 infra 已经运行了三年多，堆积了大量 technical debt。新的 infra 目标是清理历史债务，给 researcher 更好的 iteration speed。

他透露自己"其实已经不在最核心的位置了"，但认为应该做一些更重要的事情——重构 infra 就是其中之一。这与天授的经历形成了呼应：当一个系统积累了太多的不一致性和技术债务，最好的做法可能就是推倒重来，基于新的认知重新设计顶层架构。

\begin{warningbox}{Infra 的生命周期与技术债务}
翁家翌的经验表明，即使是世界顶级的 AI Infra，也有约三年的"有效生命周期"。之后技术债务累积到一定程度，patch 的成本超过重写的成本，就需要推倒重来。这与天授的情况类似——翁家翌离开后，新的维护者由于 context 不同，不可避免地引入了不一致性。\textbf{任何由人维护的系统都会逐渐腐化，关键是何时以及如何重置。}
\end{warningbox}

在 OpenAI 的 infra 体系中，researcher 不参与 infra building——他们提需求，infra 团队来实现。最终 researcher "可能到时候就改一个 flag 就好了"。

\begin{knowledgebox}{Researcher 与 Engineer 的分工}
在 OpenAI 内部，infra 和 research 有明确分工。Researcher 负责提出 idea 和需求，engineer 负责搭建正确的 infra 和快速迭代。翁家翌认为"idea is cheap"——idea 找人讨论就能出来，难的是在单位时间内能正确验证多少有效的 idea。
\end{knowledgebox}

\subsection{本章小结}

翁家翌在 OpenAI 的定位是"卖铲子的人"——不挖金矿，但给每个挖金矿的人提供最好的铲子。Post-Training RL Infra 处于技术栈最上层，生态位极高。他通过这个定位实现了 impact 的 scale up：每个使用这套 infra 训练的模型发布时，都会带上他的名字。


